\subsection{Docker}

Để đảm bảo tính nhất quán tuyệt đối của môi trường thực thi từ giai đoạn phát triển đến triển khai, đồ án lựa chọn Docker làm công cụ tiêu chuẩn để đóng gói các service. Việc áp dụng Docker trong dự án Florus không chỉ dừng lại ở việc "ảo hóa" ứng dụng mà còn tập trung vào các kỹ thuật sau:

\subsubsection{Tiêu chuẩn hóa môi trường} Mỗi thành phần trong hệ thống Florus (từ API Gateway, Backend Services đến các Data Workers) đều được định nghĩa thông qua các tệp \texttt{Dockerfile} riêng biệt. Đây được xem là tài liệu mô tả chính xác hệ điều hành nền, các thư viện phụ thuộc và biến môi trường cần thiết. Cách tiếp cận này loại bỏ hoàn toàn việc "It works on my machine", đảm bảo rằng nếu ứng dụng chạy ổn định trên máy của lập trình viên, nó cũng sẽ hoạt động chính xác trên cụm Minikube/Kubernetes.

\subsubsection{Tối ưu bằng Multi-stage Builds} 
Để tối ưu hơn trong việc sử dụng Docker, nhóm áp dụng kỹ thuật Multi-stage Builds. Thay vì gộp chung mã nguồn, công cụ biên dịch và ứng dụng vào một image khổng lồ, nhóm phân tách quy trình thành hai giai đoạn: 
\begin{itemize} 
     \item \textbf{Builder Stage:} Chứa đầy đủ các công cụ biên dịch (như Maven cho Java, GCC cho C++ hoặc Node modules đầy đủ) để thực hiện việc build mã nguồn. 
     \item \textbf{Runtime Stage:} Chỉ chứa bản binary đã biên dịch hoặc các file tĩnh cần thiết để chạy ứng dụng, hoạt động trên một nền tảng hệ điều hành rút gọn. 
\end{itemize} 

Kết quả là kích thước của các Docker Image giảm đi đáng kể (từ hàng trăm MB xuống còn vài chục MB), giúp tăng tốc độ tải image trong cluster và giảm thiểu bề mặt tấn công bảo mật do loại bỏ được các công cụ thừa.

\subsubsection{Tích hợp trong quy trình CI/CD} Trong kiến trúc tổng thể, Docker đóng vai trò là đơn vị bàn giao sản phẩm. Các Docker Image sau khi được xây dựng sẽ được đánh nhãn theo phiên bản và đẩy lên một Image Registry tập trung. Đây là tiền đề để công cụ Helm và Kubernetes có thể kéo và triển khai phiên bản ứng dụng mới nhất một cách tự động, hiện thực hóa nguyên lý "Hạ tầng bất biến".